\problemname{Village Elder}
Once upon a time there was a small village called Stackköping.
The inhabitants of Stackköping had several special traditions.
One tradition was that the oldest living villager at the end of each year had to give a New Year's speech.
Another tradition was that at most one new person could be born each year, and according to some experts, this was what eventually led to Stackköping's downfall.

During an archaeological excavation, a document was found showing the birth and death years of all $n$ people who ever lived in Stackköping.
You have obtained the document, and want to calculate how many New Year's speeches each person gave.

The New Year's speech is always the very last thing that happens each year, so no one is born or dies after the speech that occurs the same year.
If no one is alive at New Year's, then no speech is given at all.
Otherwise, a speech is always given, even if only one person is alive.

\section*{Input}

The first line contains an integer $n$ ($1 \le n \le 10^5$), the number of people.
The following $n$ lines contain two integers $f_i$ and $d_i$ ($0 \le f_i < d_i \le 10^9$): the years person number $i$ was born and died, respectively.
All values $f_i$ are distinct.

\section*{Output}
Output $n$ lines with one integer on each line, where the $i$th integer is how many New Year's speeches the $i$th person gave.

\section*{Scoring}
Your solution will be tested on a set of test groups. To earn points for a group, you must pass all test cases in that group.

\noindent
\begin{tabular}{| l | l | p{12cm} |}
  \hline
  \textbf{Group} & \textbf{Points} & \textbf{Constraints} \\ \hline
  $1$    & $30$      & $1 \le n \le 1000$ \\ \hline
  $2$    & $25$      & Everyone lived equally long, and $f_i, d_i \le 3\cdot 10^5$ \\ \hline
  $3$    & $45$      & No additional constraints.  \\ \hline
\end{tabular}
